\documentclass[reqno]{amsart}
\usepackage{amssymb, amsmath}
\usepackage[left=2.5cm, right=2.5cm, top=2.5cm, bottom=2.5cm]{geometry}
\usepackage{hyperref}
\usepackage{enumitem}

\theoremstyle{plain}
\newtheorem{theorem}{Theorem}[section]
\newtheorem{proposition}[theorem]{Proposition}
\newtheorem{lemma}[theorem]{Lemma}
\newtheorem{corollary}[theorem]{Corollary}
\theoremstyle{definition}
\newtheorem{definition}[theorem]{Definition}
\theoremstyle{remark}
\newtheorem{remark}[theorem]{Remark}

\newcommand{\Z}{\mathbb{Z}}
\newcommand{\R}{\mathbb{R}}
\newcommand{\C}{\mathbb{C}}
\newcommand{\N}{\mathbb{N}}
\newcommand{\geod}{\gamma}
\DeclareMathOperator{\limsupop}{lim\,sup}

\title{Br\"onnimann's Question~3: a group with solvable word\\ problem and irrational geodesic growth}
\author{Kirill Kondrashov}
\date{\today}

\begin{document}

\begin{abstract}
  Br\"onnimann (2016) asked whether there is a finitely generated group with solvable word
  problem whose geodesic growth series is irrational. Bodart (2023) answered this affirmatively
  with a virtually $3$-step nilpotent (\emph{Engel}) group whose geodesic growth is
  \emph{intermediate}: it satisfies $\geod(n)\asymp \exp(n^{3/5}\log n)$. This note gives a
  self-contained account of why that estimate yields a positive answer. We give complete proofs
  of the elementary bridging steps---that a sub-exponential, super-polynomial integer sequence
  cannot satisfy a linear recurrence, and that the two one-sided bounds of Bodart's estimate
  combine into a genuine growth equivalence---and we cite the deep metric-geometric inputs of
  Stoll and Pansu that underlie the growth estimate itself. A final section records that the
  argument has been machine-checked in Lean~4.
\end{abstract}

\maketitle

\section{Introduction}

Let $G$ be a group with finite generating set $S$ (closed under inverses). The
\emph{geodesic growth function} $\geod_{G,S}(n)$ counts the words of length $n$ over $S$ that are
geodesic, i.e.\ of minimal length among all words representing the same element. Its generating
series is $\Gamma_{G,S}(z)=\sum_{n\ge0}\geod_{G,S}(n)\,z^n$. Following Br\"onnimann~\cite{Bronnimann},
we call the geodesic growth \emph{rational} if $\Gamma_{G,S}$ is a rational function, and
\emph{irrational} otherwise.

\begin{quote}
\textbf{Question 3 (Br\"onnimann, 2016).} Is there a group with solvable word problem and
irrational geodesic growth?
\end{quote}

Standard families---hyperbolic groups, virtually abelian groups, and many others---have rational
geodesic growth with respect to every finite generating set, so a positive answer requires a
group whose geodesic counting sequence is genuinely non-recursive. Bodart~\cite{Bodart}
constructed such a group: a virtually $3$-step nilpotent group whose geodesic growth is
\emph{intermediate}, sitting strictly between polynomial and exponential. Intermediate growth is
incompatible with a rational series, and virtually nilpotent groups have solvable word problem;
together these answer the question.

The purpose of this note is expository and self-contained at the level of the \emph{bridging}
arguments. We isolate two elementary facts and prove them in full:
\begin{enumerate}[label=(\roman*)]
  \item an integer sequence of intermediate growth satisfies no constant-coefficient linear
    recurrence, hence has an irrational generating series (\S\ref{sec:bridge});
  \item Bodart's two one-sided bounds---an upper bound with exponent $n^{3/5}\log n$ and a lower
    bound with the smaller exponent $\kappa\,n^{3/5}\log n$, $\kappa<1$---combine into a single
    growth equivalence $\geod\asymp\exp(n^{3/5}\log n)$ (\S\ref{sec:witness}).
\end{enumerate}
The growth estimate itself (Theorem~\ref{thm:bodart}) rests on the asymptotic geometry of the word
metric on $2$-step nilpotent groups due to Stoll~\cite{Stoll} and on Pansu's theorem on the
volume growth of nilpotent groups~\cite{Pansu}; we cite these rather than reprove them.

\section{Definitions}
\label{sec:defs}

Throughout, $G=\langle \alpha\mid R\rangle$ is a group given by a presentation with generating set
$\alpha$ and relators $R$, and $S$ is a finite generating alphabet interpreted in $G$.

\begin{definition}[Solvable word problem]
The presentation $\langle\alpha\mid R\rangle$ has \emph{solvable word problem} if the relation
``the words $u$ and $v$ represent the same element of $G$'' is decidable on pairs of finite words
in $\alpha\cup\alpha^{-1}$.
\end{definition}

\begin{definition}[Geodesic growth]
A word $w\in S^\ast$ is \emph{geodesic} if no shorter word over $S$ represents the same element of
$G$. The \emph{geodesic growth function} is
\[
  \geod_{G,S}(n)=\#\{\,w\in S^{n}: w\text{ is geodesic}\,\},
\]
with generating series $\Gamma_{G,S}(z)=\sum_{n\ge0}\geod_{G,S}(n)\,z^n$.
\end{definition}

We use the classical characterization of rationality by linear recurrences (see, e.g.,
\cite[\S IV]{Bronnimann}): a power series $\sum_{n}a_n z^n$ with $a_n\in\Z$ is rational if and only
if $(a_n)$ satisfies an \emph{eventual constant-coefficient linear recurrence}, i.e.\ there exist
$d\ge1$ and $c_0,\dots,c_{d-1}\in\Z$ with
\begin{equation}\label{eq:rec}
  a_{n+d}=\sum_{i=0}^{d-1}c_i\,a_{n+i}\qquad\text{for all sufficiently large }n.
\end{equation}
Accordingly, geodesic growth is \emph{irrational} precisely when $(\geod_{G,S}(n))$ satisfies no
recurrence of the form~\eqref{eq:rec}.

Finally we recall the coarse comparison of growth functions used to state
Theorem~\ref{thm:bodart}.

\begin{definition}[Growth equivalence]
For $f,g\colon\N\to\R_{>0}$ write $f\preceq g$ if there is a constant $C\ge1$ with
$f(n)\le C\,g(Cn)$ for all large $n$, and $f\asymp g$ if $f\preceq g$ and $g\preceq f$. The
relation $\asymp$ is the standard growth equivalence of geometric group theory; note that it
permits a \emph{linear rescaling of the argument}.
\end{definition}

\section{Intermediate growth forbids a linear recurrence}
\label{sec:bridge}

We call an integer sequence $f\colon\N\to\Z$
\begin{itemize}
  \item \emph{sub-exponential} if for every $\varepsilon>0$ one has $|f(n)|\le(1+\varepsilon)^n$
    for all large $n$;
  \item \emph{super-polynomial} if for every polynomial $P$ one has $|f(n)|>P(n)$ for all large
    $n$;
  \item of \emph{intermediate growth} if it is both sub-exponential and super-polynomial.
\end{itemize}

\begin{lemma}[Recurrent sub-exponential sequences are polynomially bounded]
\label{lem:poly}
Let $f\colon\N\to\Z$ satisfy an eventual constant-coefficient linear recurrence~\eqref{eq:rec} and
be sub-exponential. Then there are $C>0$ and $r\in\N$ with $|f(n)|\le C\,(n+1)^{r}$ for all $n$.
\end{lemma}

\begin{proof}
Let $F(z)=\sum_{n\ge0}f(n)z^{n}$. Multiplying by $B(z)=1-\sum_{i=0}^{d-1}c_i z^{\,d-i}$, the
recurrence~\eqref{eq:rec} shows that all but finitely many Taylor coefficients of $B(z)F(z)$
vanish; hence $B(z)F(z)=A(z)$ for a polynomial $A$, and $F=A/B$ is a rational function.

Sub-exponentiality gives, for each $\varepsilon>0$, the bound $|f(n)|\le(1+\varepsilon)^n$
eventually, so the series $F(z)$ converges for $|z|<(1+\varepsilon)^{-1}$. As $\varepsilon>0$ is
arbitrary, $F$ is holomorphic on the open unit disk. Writing $F=A/B$ in lowest terms, every pole
of $F$ therefore lies at $|z|\ge1$. Partial fractions express
\[
  F(z)=Q(z)+\sum_{j}\frac{\alpha_j}{(1-\mu_j z)^{s_j}},
\]
with $Q$ a polynomial and, since the poles $1/\mu_j$ satisfy $|1/\mu_j|\ge1$, with $|\mu_j|\le1$.
The coefficient of $z^{n}$ in $\alpha_j(1-\mu_j z)^{-s_j}$ equals
$\alpha_j\binom{n+s_j-1}{s_j-1}\mu_j^{\,n}$, whose modulus is at most
$|\alpha_j|\,(n+1)^{s_j-1}$ because $|\mu_j|\le1$. Summing the finitely many terms (the polynomial
$Q$ contributes only to finitely many coefficients) yields $|f(n)|\le C\,(n+1)^{r}$ with
$r=\max_j s_j-1$.
\end{proof}

\begin{corollary}[Intermediate growth is irrational]
\label{cor:irr}
If $f\colon\N\to\Z$ has intermediate growth, then it satisfies no eventual
constant-coefficient linear recurrence. In particular, if $\geod_{G,S}$ has intermediate growth,
then the geodesic growth of $(G,S)$ is irrational.
\end{corollary}

\begin{proof}
Were $f$ recurrent, Lemma~\ref{lem:poly} would bound $|f(n)|$ by a polynomial, contradicting
super-polynomiality. The final clause is the definition of irrational geodesic growth.
\end{proof}

\section{Bodart's witness}
\label{sec:witness}

\subsection*{The group}
Let $a^{t}=t^{-1}at$ and $C=[a,[a,a^{t}]]$, where $[x,y]=x^{-1}y^{-1}xy$. Bodart's group is
\begin{equation}\label{eq:pres}
  G=\bigl\langle\, a,t \;\big|\; t^{2}=1,\;\; C=[a^{t},[a,a^{t}]],\;\; [C,a]=1,\;\; [C,a^{t}]=1
  \,\bigr\rangle,
\end{equation}
with geodesic alphabet $S=\{a,a^{-1},t\}$. The subgroup $\langle a,a^{t}\rangle$ is $3$-step
nilpotent of \emph{Engel} type, and $t$ acts on it as an order-two symmetry; thus $G$ is virtually
$3$-step nilpotent~\cite{Bodart}.

\begin{theorem}[Bodart {\cite[Thm.~4]{Bodart}}]
\label{thm:bodart}
For $G$ and $S$ as in~\eqref{eq:pres},
\[
  \geod_{G,S}(n)\;\asymp\;\exp\!\bigl(n^{3/5}\log n\bigr).
\]
\end{theorem}

The estimate is proved by realizing the word metric on the nilpotent part through Stoll's
asymptotic norm on $2$-step nilpotent groups~\cite{Stoll} and Pansu's volume-growth theorem for
nilpotent groups~\cite{Pansu}, refined by the geodesic combinatorics of Bishop--Elder~\cite{BishopElder}
and Bodart~\cite[\S2--\S4]{Bodart}. Concretely, Bodart establishes two one-sided bounds:
\begin{align}
  \geod_{G,S}(n)&\le c_2\,\exp\!\bigl(n^{3/5}\log n\bigr) &&\text{eventually,}\label{eq:upper}\\
  \geod_{G,S}(n)&\ge c_1\,\exp\!\bigl(\kappa\, n^{3/5}\log n\bigr) &&\text{eventually,}\label{eq:lower}
\end{align}
for constants $c_1,c_2>0$ and some $\kappa\in(0,1]$. The upper bound~\eqref{eq:upper} carries the
exponent constant $1$; the lower bound~\eqref{eq:lower} only controls the exponent up to a
\emph{smaller} constant $\kappa<1$ coming from the explicit geodesic family of \cite[\S4]{Bodart}.

\begin{remark}
The two exponent constants \emph{cannot} be equated: since
$\exp(\kappa\,n^{3/5}\log n)/\exp(n^{3/5}\log n)\to0$ super-polynomially, no constant $c_1'>0$
makes $c_1'\exp(n^{3/5}\log n)\le\geod_{G,S}(n)$ hold. It is the argument-rescaling freedom built
into $\asymp$ that reconciles them, as the next lemma shows.
\end{remark}

\begin{lemma}[Reconciliation of exponent constants]
\label{lem:coarse}
Set $M(n)=\exp(n^{3/5}\log n)$ and $M_\kappa(n)=\exp(\kappa\, n^{3/5}\log n)$ with $\kappa\in(0,1]$.
If \eqref{eq:upper} and \eqref{eq:lower} hold, then $\geod_{G,S}\asymp M$.
\end{lemma}

\begin{proof}
The upper bound~\eqref{eq:upper} gives $\geod_{G,S}\preceq M$ directly (take $C=\max(c_2,1)$,
argument scale $1$). For $M\preceq \geod_{G,S}$, choose an integer $k\ge\kappa^{-5/3}$; then
$\kappa\,k^{3/5}\ge1$. For $n\ge1$, using $\log(kn)=\log k+\log n\ge\log n\ge0$,
\[
  M_\kappa(kn)=\exp\!\bigl(\kappa\,k^{3/5}\,n^{3/5}\,\log(kn)\bigr)
  \ge\exp\!\bigl(n^{3/5}\log n\bigr)=M(n).
\]
Combining with~\eqref{eq:lower} evaluated at $kn$ gives, for all large $n$,
\[
  M(n)\le M_\kappa(kn)\le c_1^{-1}\,\geod_{G,S}(kn).
\]
Thus $M(n)\le C\,\geod_{G,S}(Cn)$ with $C=\max(k,c_1^{-1})$, i.e.\ $M\preceq\geod_{G,S}$.
\end{proof}

\begin{corollary}
\label{cor:inter}
$\geod_{G,S}$ has intermediate growth.
\end{corollary}

\begin{proof}
For every $\varepsilon>0$ one has $M(n)=\exp(n^{3/5}\log n)<(1+\varepsilon)^n$ for large $n$
(since $n^{3/5}\log n=o(n)$), so by~\eqref{eq:upper} $\geod_{G,S}$ is sub-exponential. For every
polynomial $P$ one has $P(n)<\exp(\kappa\,n^{3/5}\log n)=M_\kappa(n)$ for large $n$ (since
$\kappa\,n^{3/5}\log n$ eventually exceeds $\deg(P)\log n+O(1)$), so by~\eqref{eq:lower}
$\geod_{G,S}$ is super-polynomial.
\end{proof}

\section{Solvable word problem}
\label{sec:wp}

\begin{lemma}
\label{lem:wp}
Finitely generated virtually nilpotent groups have solvable word problem; in particular so does
the group $G$ of~\eqref{eq:pres}.
\end{lemma}

\begin{proof}
A finitely generated nilpotent group is linear over $\Z$ and hence residually finite
(Mal'cev~\cite{Malcev}); residual finiteness passes to finite-index overgroups, so $G$ is
residually finite. A finitely presented residually finite group has solvable word problem: to
decide whether a word $w$ is trivial, run in parallel an enumeration of consequences of the
relators (which halts if $w=1$) and an enumeration of finite quotients (which exhibits a quotient
separating $w$ from $1$ if $w\ne1$). Constructively, $G$ embeds its nilpotent finite-index
subgroup into an explicit group of integer coordinate tuples on which multiplication and equality
are computable, and equality in $G$ is decided by comparing normal forms in that model.
\end{proof}

\section{Main result}

\begin{theorem}
\label{thm:main}
The group $G$ of~\eqref{eq:pres}, with generating alphabet $S=\{a,a^{-1},t\}$, has solvable word
problem and irrational geodesic growth. Consequently Br\"onnimann's Question~3 has a positive
answer.
\end{theorem}

\begin{proof}
Solvability of the word problem is Lemma~\ref{lem:wp}. By Theorem~\ref{thm:bodart} and
Lemma~\ref{lem:coarse} the geodesic growth satisfies $\geod_{G,S}\asymp\exp(n^{3/5}\log n)$; by
Corollary~\ref{cor:inter} it is of intermediate growth, and by Corollary~\ref{cor:irr} it is
therefore irrational. Thus $(G,S)$ realizes both properties demanded by Question~3.
\end{proof}

\section{Formalization in Lean}
\label{sec:lean}

The argument of this note has been formalized and machine-checked in Lean~4 with
\texttt{mathlib}. Fully kernel-checked from the base axioms ($\mathtt{propext}$,
$\mathtt{Quot.sound}$, $\mathtt{Classical.choice}$) are: the reduction of Question~3 to the
existence of a witness (Theorem~\ref{thm:main}'s packaging); the recurrence-to-polynomial bridge
(Lemma~\ref{lem:poly}) and its irrationality corollary (Corollary~\ref{cor:irr}); the
growth-equivalence reconciliation (Lemma~\ref{lem:coarse}); the coordinate model of $G$ with
verification of the four relators of~\eqref{eq:pres}; and the resulting word-problem decision
procedure (Lemma~\ref{lem:wp}). The two deep analytic halves of Theorem~\ref{thm:bodart}---the
upper bound~\eqref{eq:upper} via Stoll/Pansu and the lower bound~\eqref{eq:lower} via Bodart's
\S4 geodesic family---are currently recorded as two named axioms, isolating exactly the classical
metric-geometric inputs that remain to be internalized.

\medskip
\noindent Source, build instructions, and the full axiom report:
\url{https://github.com/kirill-kondrashov/lean-misc} (see the \emph{Br\"onnimann Question~3}
section of the repository \texttt{README}).

\begin{thebibliography}{9}

\bibitem{Bronnimann}
J.~M. Br\"onnimann, \textit{Geodesic Growth of Groups}, PhD thesis, Universit\'e de Neuch\^atel,
2016.

\bibitem{Bodart}
C.~Bodart, \textit{Intermediate geodesic growth in virtually nilpotent groups},
arXiv:2306.10381v3, 2025.

\bibitem{BishopElder}
A.~Bishop and M.~Elder, \textit{A virtually $2$-step nilpotent group with polynomial geodesic
growth}, Algebra Discrete Math. \textbf{33} (2022), no.~1, 21--28.

\bibitem{Stoll}
M.~Stoll, \textit{On the asymptotics of the growth of $2$-step nilpotent groups},
J.~London Math. Soc. (2) \textbf{58} (1998), no.~1, 38--48.

\bibitem{Pansu}
P.~Pansu, \textit{Croissance des boules et des g\'eod\'esiques ferm\'ees dans les
nilvari\'et\'es}, Ergodic Theory Dynam. Systems \textbf{3} (1983), no.~3, 415--445.

\bibitem{Malcev}
A.~I. Mal'cev, \textit{On homomorphisms onto finite groups}, Ivanov. Gos. Ped. Inst. Uchen. Zap.
\textbf{18} (1958), 49--60; English transl.\ Amer. Math. Soc. Transl. (2) \textbf{119} (1983),
67--79.

\end{thebibliography}

\end{document}
